\section{Introduction}

\subsection{Background and Motivation}

The digitization of healthcare systems over the past two decades has led to exponential growth in clinical data generated across care settings. Despite widespread adoption of Electronic Health Record (EHR) systems, a substantial proportion of clinically valuable information remains locked in unstructured formats, including physician dictations, bedside observations, referral notes, and patient-provider conversations. Studies have consistently demonstrated that such data, when appropriately processed, contributes meaningfully to clinical tasks such as phenotyping and quality measurement, yet it remains largely inaccessible to automated systems in its raw form \cite{hong2020nlp2fhir}. Structured data, stored in coded fields and controlled vocabularies, enables systematic querying and analysis; unstructured data, by contrast, often captures the richness, nuance, and context of clinical encounters that structured forms cannot convey \cite{hong2017integrating, brens2025semantic}. Across 172 clinical phenotypes derived from MIMIC-III, textual features outperform structured measurements alone for 84\% of phenotypes, with their combination yielding statistically significant improvements where clinical notes capture context not reflected in tabular measurements \cite{moldwin2021phenotyping}. Bridging this divide has long been recognized as a fundamental informatics challenge: extracting meaningful, computable representations from free-text and audio-based clinical content requires not only language understanding but also domain-specific knowledge and semantic reasoning \cite{builtjes2025leveraging}.

This documentation burden carries direct consequences for clinicians and patients alike. Primary care physicians spend a disproportionate share of their working time on administrative documentation rather than direct patient care, a pattern documented consistently across healthcare systems internationally \cite{saadat2025enhancing}. Documentation workload is a well-established driver of clinician burnout and professional attrition \cite{klusty2024clinical}, and documentation delays introduce latency between clinical events and their representation in the EHR, undermining care team access to timely information and increasing transcription error risk. Clinical audio compounds this problem: patient-clinician encounters, whether in outpatient consultations, emergency settings, or telehealth platforms, generate continuous streams of spoken dialogue with diagnostically relevant content, yet converting this audio into structured, actionable data remains overwhelmingly manual, relying on transcriptionists or clinicians themselves to document encounters post hoc \cite{saadat2025enhancing, klusty2024clinical}. Automating this pathway addresses a systemic inefficiency with measurable consequences for workforce sustainability, patient safety, and care quality. Reducing the manual overhead of clinical documentation while preserving its fidelity and interoperability is the central aim of this work.

\subsection{The Role of HL7 FHIR in Healthcare Interoperability}

HL7 Fast Healthcare Interoperability Resources (FHIR) has emerged as the leading international standard for the representation and exchange of healthcare information. By defining a modular, resource-based data model encompassing clinical entities such as \texttt{Patient}, \texttt{Encounter}, \texttt{Observation}, \texttt{Condition}, and \texttt{MedicationRequest}, FHIR enables semantic interoperability across disparate health information systems \cite{hl7fhir2022r4b}. Recent work applying large language models to FHIR transformation has demonstrated a broad range of applications, from clinical decision support and patient data exchange to population health management and regulatory reporting, while also highlighting persistent implementation challenges such as terminology binding, resource profiling complexity, and inconsistent adoption across jurisdictions \cite{li2024fhirgpt}. Mandated or preferred by major healthcare regulators, including the U.S.\ Office of the National Coordinator for Health IT, FHIR has seen accelerating global uptake. Despite this progress, the vast majority of FHIR-compliant data structures are populated from existing structured EHR fields; the generation of FHIR resources from unstructured sources, particularly spoken clinical audio, remains largely unsolved. Automating this pathway would enable real-time population of clinical records and support downstream analytics, care coordination, and research reuse.

\subsection{Automatic Speech Recognition in Healthcare}

Automatic Speech Recognition (ASR) has advanced considerably in recent years, driven by deep learning architectures and large-scale pretraining. Radford et al.\ introduced Whisper, a general-purpose ASR model trained on 680,000 hours of multilingual, weakly supervised audio data, demonstrating near-human transcription accuracy and robust zero-shot generalization across diverse acoustic domains \cite{radford2023whisper}. The accessibility of such pretrained models makes them compelling candidates for clinical audio processing pipelines, eliminating the need to train ASR systems from scratch.

The application of ASR in healthcare carries distinctive challenges. Medical terminology is acoustically dense with domain-specific vocabulary, such as drug names, anatomical terms, and diagnostic codes, and is often produced under difficult recording conditions: background noise, overlapping speakers, and informal conversational register. Independent evaluation on real patient-clinician conversations has confirmed that, even under controlled conditions, clinical dialogue yields word error rates of 8.8--10.5\% alongside variable speaker diarization accuracy \cite{tran2022asr}. The PriMock57 dataset, comprising 57 recorded mock primary care consultations with manual transcriptions, provides a key benchmark for evaluating ASR systems in this setting \cite{papadopoulos2022primock57}. Adedeji et al.\ showed that applying LLMs, particularly through Chain-of-Thought (CoT) prompting, to refine ASR-generated transcripts improved both general and medical-concept word error rates, as well as diarization accuracy and semantic coherence \cite{adedeji2024sound}. These findings suggest that a hybrid pipeline combining a pretrained ASR system with LLM-based post-processing is a practical approach to medical transcription without requiring custom model development.

\subsection{From Transcription to Structured Clinical Data}

Transcription alone is insufficient for clinical informatics applications. To be computationally actionable, spoken clinical content must be transformed into structured representations grounded in standard terminologies and data models, requiring NLP pipelines capable of named entity recognition, relation extraction, negation detection, and normalization against controlled vocabularies such as SNOMED CT, LOINC, and RxNorm. Extraction of clinical information from unstructured data, including clinical notes and transcribed encounters, has been demonstrated as feasible using open-source LLMs in zero-shot settings, though accuracy varies substantially by entity type and documentation style \cite{builtjes2025leveraging}. Mapping extracted concepts to FHIR resources introduces further complexity: each resource type carries specific structural requirements, terminology bindings, and cardinality constraints that must be satisfied to produce conformant output. The challenge of designing a pipeline that traverses the full pathway from raw audio to validated FHIR resource instances, integrating ASR, NLP, terminology normalization, and resource generation, has not been comprehensively addressed in the literature; existing approaches tend to address individual stages in isolation \cite{hong2020nlp2fhir, riquelme2025llm}.

\subsection{Problem Statement}

Despite the individual maturity of its constituent technologies (ASR, clinical NLP, and FHIR modelling), no established solution integrates these into a coherent, end-to-end pipeline that transforms raw clinical audio into validated, interoperable FHIR resources. Existing approaches address individual stages in isolation: ASR systems are evaluated for transcription accuracy without downstream clinical structuring; NLP pipelines operate on pre-existing text rather than live transcriptions; and FHIR generation tools assume structured or semi-structured input rather than conversational speech. The full pathway from an audio recording of a clinical encounter to a computable, standards-conformant clinical record thus remains an open engineering problem \cite{hong2020nlp2fhir, riquelme2025llm}.

This study therefore poses the following question: can an automated pipeline built from existing pretrained components reliably transcribe primary care consultations and produce clinically meaningful, FHIR R4B-conformant resources, without requiring custom model development or manual post-editing?

\subsection{Research Objectives and Scope}

This study presents the design and evaluation of an automated pipeline for mapping unstructured clinical audio to HL7 FHIR Release 4B-compliant data structures \cite{hl7fhir2022r4b}. Rather than developing novel ASR or NLP models from scratch, the proposed system integrates and orchestrates existing pretrained components, including Whisper \cite{radford2023whisper} for transcription and LLM-based processing for clinical entity identification \cite{adedeji2024sound}, within a unified, modular architecture. The pipeline is developed and evaluated using the PriMock57 dataset of primary care mock consultations \cite{papadopoulos2022primock57}; using mock consultations with professional actors minimizes ethical risks associated with patient-identifiable data while maintaining clinical realism. The novelty of this work lies in the pipeline design: the systematic integration of these components, the application of LLM-driven post-processing for both transcription refinement and clinical information extraction, and the mapping of extracted entities to validated FHIR resource structures.

The remainder of this paper is organized as follows. Section~2 reviews related work spanning clinical ASR, medical NLP, and FHIR-based data modeling. Section~3 describes the proposed pipeline architecture and its constituent components. Section~4 details the experimental setup and evaluation methodology. Section~5 presents results and discussion, and Section~6 concludes with implications for clinical practice and directions for future research.
